\documentclass[10pt,a4paper]{amsart}
\usepackage[margin=19mm]{geometry}
\usepackage{amsmath,amssymb,mathtools}
\usepackage[scr=rsfs,cal=euler]{mathalfa}
\usepackage{xfrac}
\usepackage[unicode,hidelinks,bookmarks=false]{hyperref}
\newcommand{\ie}{i.e.\ }
\newcommand{\cf}{cf.\ }
\newcommand{\resp}{respectively\ }
\newcommand{\Subsection}{\subsection}
\providecommand{\nobreakdash}{\nobreak\hbox{-}}
\setlength{\emergencystretch}{2em}
\multlinegap0pt
\usepackage{bm}
\usepackage{bbm}
\usepackage{dsfont}
\usepackage{xspace}
\usepackage{cancel}
\usepackage{mathtools}
\usepackage{amsthm}
\makeatletter
\@ifundefined{theorem}{\newtheorem{theorem}{Theorem}}{}
\@ifundefined{lemma}{\newtheorem{lemma}[theorem]{Lemma}}{}
\makeatother
\newcommand{\eps}{\varepsilon}
\renewcommand{\geq}{\geqslant}
\renewcommand{\leq}{\leqslant}
\title[Class numbers and integer points on some Pellian surfaces]{Class numbers and integer points on some Pellian surfaces}
\author{Yijie Diao}
\date{Source: arXiv:2408.03774v2 (CC BY 4.0)}
\begin{document}
\maketitle

\noindent\emph{Source and licence.} Class numbers and integer points on some Pellian surfaces. Version arXiv:2408.03774v2 (CC BY 4.0), distributed under the Creative Commons Attribution 4.0 International licence, as recorded in the arXiv licence metadata for this exact version and shown on the abstract page https://arxiv.org/abs/2408.03774v2. Full rights notice: licenses/arxiv-2408.03774-CC-BY-4.0.txt.\par\smallskip

\noindent\textit{Editorial scope.} Counted unit: the Lemma whose source label is \texttt{Reu+FJ} in the article (source0.tex lines 330--333, source label \texttt{Reu+FJ}), delivered together with the article's own complete original proof of it (source0.tex lines 334--377, one \texttt{proof} environment of 44 lines). No mathematics is written, repaired or re-derived: statement and proof are the article's own text line for line. Required context retained uncounted: thm-Reu (lines 317--319); thm-FJ (lines 325--327). Every definition, display and label of those lines is retained unchanged. Omitted from this package and not redistributed: 1--316, 320--324, 328--329, 378--810. Nothing inside a retained excerpt is omitted or altered: every retained body is delivered exactly as the article's lines are written, so no omission receipt of any kind accompanies this package. Citations: the retained excerpts carry no citation command at all, so no bibliography entry of the article is transcribed and no reference is redistributed. Reference adaptation: none was needed and none was made; every reference inside the delivered unit resolves inside it, so no unresolved reference or citation is produced. Printed slips kept literal: no printed word, symbol, hypothesis or number of the counted statement or of its proof was altered. Layout and class: the article's own class is \texttt{amsart} with packages including amsfonts, amsthm, amsmath, amssymb, geometry, mathrsfs, enumitem, hyperref; the wrapper is the lane's pinned \texttt{amsart} editorial wrapper at 10pt on A4, which restates the theorem-like environments the retained text needs and adds the article's own macro definitions that the retained text needs (\texttt{\textbackslash eps}, \texttt{\textbackslash geq}, \texttt{\textbackslash leq}), with extra wrapper packages (bm, bbm, dsfont, xspace, cancel, mathtools). The delivered numbering is the unit's own. Assets: no retained span contains a figure environment, a graphics-inclusion command, a PDF-page inclusion, a table or a TikZ picture; no image file is copied into this package, the package declares no asset dependency and no image file is redistributed with it. Licence: version arXiv:2408.03774v2 of this article is distributed under the Creative Commons Attribution 4.0 International licence, as recorded in the arXiv licence metadata for this exact version and shown on the abstract page https://arxiv.org/abs/2408.03774v2. A full rights notice is delivered at licenses/arxiv-2408.03774-CC-BY-4.0.txt. Characters: every retained excerpt is pure ASCII, so the delivered engine, XeLaTeX with its default Latin Modern text font, reports no missing character.
\begin{align}\label{thm-Reu}
	N_\theta (D_1, D_2, U_1, U_2) \ll (D_2 U_2)^\eps \min \big((U_1U_2M)^\frac{1}{2} + U_1 + U_2, (D_1D_2M)^\frac{1}{2} + D_1 + D_2 \big), 
\end{align}

\begin{eqnarray}\label{thm-FJ}
	N_\theta (D_1, D_2, U_1, U_2) \ll (D_1D_2)^\eps \big((D_1D_2)^\frac{1}{2} + D_1U_2\big).
\end{eqnarray}

\begin{lemma}\label{Reu+FJ}
	Let $\epsilon > 0$. Assume that $D_1 \leq D_2$, $D_1 D_2 \leq B$ and $D_1U_1^2 = D_2 U_2^2 = B$, then for any $\theta = \pm 1$ or $\pm 2$, we have
	$$N_\theta(D_1, D_2, U_1, U_2) \ll B^{\frac{7}{12} + \epsilon}.$$
\end{lemma}

\begin{proof}
	When $D_1 U_2 \leq B^{\frac{7}{12}}$, then (\ref{thm-FJ}) implies 
	$$N_\theta(D_1, D_2, U_1, U_2) \ll B^\epsilon (B^{\frac{1}{2}} + B^{\frac{7}{12}}) \ll B^{\frac{7}{12} + \epsilon}.$$
	
    Next we observe that if $\log (D_1D_2) = k \log B$, then 
    \begin{eqnarray}\label{log-M}
    	m: = \frac{\log M}{\log B} = \frac{9}{16} k (2 - k).
    \end{eqnarray}
    Since $0 \leq k \leq 1$, the function $m(k)$ is an increasing function. 
	
	When $D_1 U_2 > B^{\frac{7}{12}}$, then considering $D_2U_2^2 = B$, we know 
	$$D_1 D_2^{-\frac{1}{2}} = D_1U_2B^{-\frac{1}{2}} > B^{\frac{1}{12}}.$$ 
	Since $D_1D_2 \leq B$, we obtain
	$$D_2D_1^{-1} = (D_1D_2)^{\frac{1}{3}} \cdot (D_1D_2^{-\frac{1}{2}})^{-\frac{4}{3}} < B^{\frac{1}{3} - \frac{1}{9}} =    B^{\frac{2}{9}}.$$
	On the other hand, using $D_1 \leq D_2$, we know  $$D_1D_2 = (D_1^{-1}D_2)^3 \cdot (D_1D_2^{-\frac{1}{2}})^4 \geq (D_1 D_2^{-\frac{1}{2}})^4 > B^{\frac{1}{3}}.$$ 
	In other words, we have $k > \frac{1}{3}$. Since $m(k)$ is increasing when $0 \leq k \leq 1$, we have $m(k) > m(\frac{1}{3}) = \frac{5}{16}$. Hence we have
	\begin{eqnarray}\label{estimate-M}
		M > B^{\frac{5}{16}}.
	\end{eqnarray}
	Therefore, it follows that 
	$$M > B^{\frac{2}{9}} \geq D_2D_1^{-1}, \text{ which implies } (D_1D_2M)^{\frac{1}{2}} > D_2 \geq D_1.$$
	If we further assume that $k \leq \frac{2}{3}$, then by (\ref{thm-Reu}) we have
	$$N_\theta(D_1, D_2, U_1, U_2) \ll B^\epsilon (D_1D_2M)^{\frac{1}{2}} \ll B^{\frac{1}{2}(\frac{2}{3} + m(\frac{2}{3})) + \epsilon} = B^{\frac{7}{12} + \epsilon},$$
	since $m(k) + k$ is also an increasing function. 
	
The remaining case is when $D_1 U_2 > B^{\frac{7}{12}}$ and $k > \frac{2}{3}$. Note that $$U_1U_2 = (BD_1^{-1})^{\frac{1}{2}} \cdot (BD_2^{-1})^{\frac{1}{2}} = B^{1 - \frac{k}{2}}.$$ 
By (\ref{estimate-M}), we obtain that
$$M > B^{\frac{5}{16}} > B^{\frac{1}{9}} \geq (D_2D_1^{-1})^{\frac{1}{2}} = U_1U_2^{-1}, \text{ which implies }  (U_1U_2M)^{\frac{1}{2}} > U_1 \geq U_2.$$ 
We also have
$$(U_1U_2M)^{\frac{1}{2}} = B^{\frac{1}{2}(1  - \frac{k}{2} + m(k))}.$$
Note that the function 
$- \frac{k}{2} + m(k) = \frac{1}{16}k(10-9k)$
 is decreasing when $k > \frac{5}{9}$. Therefore, by the assumption $k > \frac{2}{3}$ and (\ref{thm-Reu}) we have
$$N_\theta(D_1, D_2, U_1, U_2) \ll B^\epsilon (U_1U_2M)^{\frac{1}{2}} \ll B^{\frac{1}{2}(1 - \frac{1}{3} + m(\frac{2}{3})) + \epsilon} = B^{\frac{7}{12} + \epsilon}.$$
	
	%When $D_1M > D_2$, we have $(D_1D_2M)^{\frac{1}{2}} > D_2 \geq D_1$. 
	
	%When $D_2 \leq B^{\frac{7}{12}}$, we have
	
	%When $D_1 U_2 > B^{\frac{7}{12}}$, we have $D_2 D_1^{-1} < B^{\frac{5}{12}}$ since $D_2 U_2 \leq D_2U_2^2 = B$. 
	
	%When $D_1 U_2 > B^{\frac{7}{12}}$ and $D_1D_2 \leq B^{\frac{2}{3}}$, then by $(\ref{thm-Reu})$, we have
	%$$N_\eta(D_1, D_2, U_1, U_2) \ll B^\epsilon (D_1D_2M)^{\frac{1}{2}} \ll .$$
\end{proof}

\end{document}
